% !PRINTED 272
\weberpage{272}
%<<CONT
for which one obtains from (2) the corresponding permutations:
\weq{12}{\begin{array}{lll}
\omega, & A, & B,\\
-\dfrac{1}{\omega}, & \varrho^2 A, & \varrho B,\\
\omega+1, & e^{\frac{(n+1)\pi i}{6}}\sigma^2 A, & e^{-\frac{(n+1)\pi i}{6}}\sigma B,
\end{array}}
and from this one derives in the same way the equations:
\weq{13}{\begin{aligned}
n&=3, & A&=0,\\
n&=11, & A&=1,\\
n&=19, & A^5-7A^2-B&=0.
\end{aligned}}
3. If $n \equiv 1$ (mod 4), one can proceed in the same way with the functions
\weq{14}{\begin{aligned}
8A &= u^4v^4 - u_1^4v_1^4 - u_2^4v_2^4,\\
B &= u^4v^4u_1^4v_1^4 + u^4v^4u_2^4v_2^4 - u_1^4v_1^4u_2^4v_2^4,
\end{aligned}}
for which one obtains the corresponding permutations:
\weq{15}{\begin{array}{lll}
\omega, & A, & B,\\
-\dfrac{1}{\omega}, & \varrho\,A, & \varrho^2 B,\\
\omega+1, & e^{\frac{(n+1)\pi i}{3}}\sigma\,A, & e^{-\frac{(n+1)\pi i}{3}}\sigma^2 B.
\end{array}}
Only the first case $n=5$ leads here to a simple result:
\weq{16}{n=5, \quad A=1.}
For $n=5$ a still simpler form of the transformation equation can be obtained.

For if we apply to the three functions
\weq{17}{w = \frac{u_1^2v_2^2+u_2^2v_1^2}{uv}, \quad w_1 = \frac{u_2^2v^2-v_2^2u^2}{u_1v_1}, \quad w_2 = \frac{u_1^2v^2-v_1^2u}{u_2v_2}} % [sic] German source prints u without the square in the w_2 numerator.
the substitutions $\bigl(\omega, -\frac{1}{\omega}\bigr)$, $(\omega, \omega+1)$, then under the assumption $n \equiv 5$ (mod 8) the corresponding permutations result:
\[
\begin{array}{llll}
\omega, & w, & w_1, & w_2,\\
-\dfrac{1}{\omega}, & w, & w_2, & w_1,\\
\omega+1, & e^{-\frac{\pi i}{24}(n-5)}w_1, & e^{-\frac{\pi i}{24}(n-5)}w, & e^{\frac{\pi i}{12}(n-5)}w_2.
\end{array}
\]
